\subsection{Digital Certificates}\label{sec:binding_keys_to_parties:digital_certificate}

All trust models encode and validate trust relationships through some sort of digital certificates. In general, a \textbf{digital certificate} (or simply certificate) is an electronic document cryptographically binding a key pair with a unique identity (or pseudo-identity) of a party (e.g., users, services and agents in general). A certificate usually contains the public key, the unique identity of the party, further attributes of the party, and a digital signature (\Cref{sec:digital_signatures}) of the certificate itself. Such a digital signature may have been created with the private key corresponding to the public key contained in the certificate --- in which case, the certificate is said to be \textbf{self-signed} --- or with the private key of a (often widely trusted or authoritative) third party typically called \textbf{\gls{ca}}. A certificate usually has a start and an end time for its validity --- often (but not necessarily) corresponding to the cryptoperiod of the key pair (\Cref{sec:key_management:lifecycle:operational}) but can even be revoked in advance (\Cref{sec:binding_keys_to_parties:revocation}).

\warningBlock{Verifying possession of private keys}{A certificate binding a party to a public key allows the party to claim control of the corresponding private key. Therefore, the procedure for releasing a certificate needs to include a verification step, that is, ensure that the recipient party actually controls the private key. Obviously, the lack of such a verification step invalidates any authenticity claims.}

\subsubsection{Validation of Digital Certificates}\label{sec:binding_keys_to_parties:digital_certificate:validation}

The validation of a certificate is a complex procedure, but roughly involves (at least) the following steps:

\begin{enumerate}
	\item check the digital signature of the certificate;
	\item compare the validity start and end time of the certificate against the current time;
	\item ensure the certificate has not been revoked --- this step may vary depending on the revocation procedure employed (\Cref{sec:binding_keys_to_parties:revocation});
	\item establish a transitive trust path linking the certificate to a predefined \textbf{root of trust} --- this step surely varies depending on the chosen trust model. This trust path is often called \textbf{trust chain} (especially in \gls{pki});
	\item consider policy and application-level checks (e.g., the certificate may be valid but not properly formatted, the policy does not accept it due to the use of insecure chosen ciphers, the attributes contained in the certificate do not match the intended usage).
\end{enumerate}	

\remarkBlock{More on the establishment of trust paths}{One of the steps to validate a certificate is establishing a trust path: this step is not only cryptographic but also concerns trust assumptions and trust endorsements. In fact, a party who wants to validate a certificate (called \textbf{validator}) ultimately needs to have assumed trust on other parties in the real world --- that is, having defined one or more roots of trust. \textbf{Without any prior real-world trust assumptions, is it impossible to have trust}. A root of trust may either be another party (e.g., a user) or an organization (i.e., \glspl{ca}), and the definition of a root of trust varies depending on the underlying trust model. Also, each party may have its own roots of trust. Consequently, the establishment of a trust path is the (often recursive or iterative) procedure through which the validator ensures that the certificate being validate is (either directly or indirectly) endorsed by one or more (or enough) roots of trust.}

\warningBlock{Misinterpreting certificates}{A valid certificate does not inherently imply that the owner of the certificate is trusted, but only that some root of trust endorsed the key pair and the attributes contained therein. Moreover, these attributes are often misinterpreted, yielding to misplace trust. For instance, assuming that a certificate for a website guarantees the honesty and reliability of the website owner is a common misconception. In fact, scam websites often exhibit certificates whose attributes (e.g., website name and url) resemble those of the original websites (a practice called \textbf{domain name mismatch}). In this sense, an \texttt{https} connection only guarantees the security of the communication, not the authenticity or trustworthiness of the website owner.}
